\section{Travaux connexes}\label{sec-related}

\paragraph{Détection d'objets en une passe.}
La famille YOLO~\cite{redmon2016yolo} a imposé le paradigme de la
détection en une seule passe : l'image entière traverse une fois le
réseau, qui prédit conjointement positions et classes, ce qui permet
l'inférence temps réel indispensable à la vidéoprotection. Les versions
successives ont notamment enrichi les augmentations de données :
l'augmentation \emph{mosaic}, introduite avec
YOLOv4~\cite{bochkovskiy2020yolov4}, assemble quatre images en une et
expose le réseau à des objets réduits et à des contextes variés.
YOLOv8~\cite{yolov8_ultralytics} en est l'implémentation de référence
actuelle chez Ultralytics ; sa déclinaison nano, la plus légère, est
celle utilisée ici.

\paragraph{Le cas des petits objets.}
La difficulté des détecteurs sur les petits objets est documentée depuis
la création du benchmark COCO~\cite{lin2014coco}, au point que son
protocole officiel ventile les métriques en trois catégories de taille
(petit, moyen, grand), convention que nous reprenons
(\cref{sec-protocole}). La cause est structurelle : après les
sous-échantillonnages successifs du réseau, un véhicule de moins de
$32 \times 32$ pixels ne pèse plus que quelques activations, souvent
insuffisantes pour déclencher une détection. C'est précisément la
configuration des véhicules éloignés ou masqués dans les embouteillages
qui nous occupe.

\paragraph{Jeux de données de trafic et écart de domaine.}
Les benchmarks historiques de surveillance de trafic, comme
UA-DETRAC~\cite{wen2020uadetrac} (100 séquences filmées en Chine),
proviennent de contextes routiers peu représentatifs des métropoles
denses des pays en développement. BMD-45~\cite{sharma2026bmd45} comble
ce manque avec 45\,000 images issues de plus de 3\,600 caméras CCTV
opérationnelles à Bengaluru, annotées en 14 catégories fines incluant
des modes de transport régionaux. Ses auteurs quantifient l'écart de
domaine : des détecteurs entraînés sur UA-DETRAC atteignent $33{,}6\,\%$
de mAP@0,5:0,95 sur leur jeu, contre $83{,}8\,\%$ pour des modèles
entraînés en domaine, un facteur $2{,}5$ qui justifie à lui seul la
démarche d'adaptation entreprise ici.

\paragraph{Vision par ordinateur pour le trafic en Afrique.}
Les travaux africains sur le sujet restent rares.
Mugizi \emph{et al.}~\cite{mugizi2026traffic} ont développé en Ouganda
un système complet de surveillance du trafic (détection de véhicules,
lecture de plaques et estimation de vitesse) atteignant $97{,}9\,\%$ de
mAP en détection de plaques. Leur travail couvre le pipeline de bout en
bout ; le nôtre est complémentaire : nous isolons le premier maillon, la
détection, et son point de fragilité documenté, les petits objets. À
notre connaissance, il n'existe aucun jeu de données public de vision
par ordinateur consacré au trafic béninois, ce qui motive le choix d'un
domaine structurellement comparable (\cref{sec-donnees}) et constitue la
première limite discutée en \cref{sec-discussion}.
